\documentclass[10pt, letterpaper]{article}

% Inhaltsverzeichnis für Pakettypen (nur für Übersicht im Header, wird nicht im Dokument angezeigt)
% 1. Seitenlayout und Ränder
% 2. Sprache und Zeichensatz
% 3. Mathematik und Theorem-Umgebungen
% 4. Eigene Makros
% 5. Diagramme und Grafiken
% 6. Tabellen und Aufzählungen
% 7. Inhaltsverzeichnis
% 8. Abschnittsüberschriften
% 9. Abstrakt-Umgebung
% 10. Todos/Notizen
% 11. Rahmen/Box-Umgebungen
% 12. Python-Integration
% 13. Literaturverwaltung
% 14. Hyperlinks
% 15. Absatzeinstellungen
% 16. Umgebungen
% 17  Graphik
% 18  Extra
% 00. Titel und Autor

% --- 1. Seitenlayout und Ränder ---
\usepackage[margin=3cm]{geometry}

% --- 2. Sprache und Zeichensatz ---
\usepackage[english]{babel}
\usepackage[T1]{fontenc}
\usepackage[utf8]{inputenc}

% --- 3. Mathematik und Theorem-Umgebungen ---
\usepackage{amsmath, amssymb, amsthm}
\usepackage{mathrsfs}
\DeclareMathOperator{\WF}{WF}

% --- 4. Eigene Makros ---
\usepackage{xcolor}
\newcommand{\SKP}{\langle\cdot,\cdot\rangle}
\newcommand{\R}{\mathbb{R}}
\newcommand{\N}{\mathbb{N}}
\newcommand{\Q}{\mathbb{Q}}
\newcommand{\Z}{\mathbb{Z}}
\newcommand{\C}{\mathbb{C}}
\newcommand{\entwurf}[1]{\textcolor{red}{#1}}

% --- 5. Diagramme und Grafiken ---
\usepackage{graphicx}
\graphicspath{{../../Library/Images/}}
\usepackage[export]{adjustbox}
\usepackage{tikz}
\usetikzlibrary{decorations.pathreplacing, arrows.meta, positioning}
\usepackage{tikz-cd}

% --- 6. Tabellen und Aufzählungen ---
\usepackage{enumitem}
\setlist[itemize]{left=0.5cm}

\newenvironment{romanenum}[1][]
  {%
    \ifx&#1&
    \else
      \textbf{#1}\quad
    \fi
    \begin{enumerate}[label=\roman*)]
  }
  {%
    \end{enumerate}%
  }

% --- 7. Inhaltsverzeichnis ---
\usepackage{tocloft}
\renewcommand{\cftsecfont}{\footnotesize}
\renewcommand{\cftsubsecfont}{\footnotesize}
\renewcommand{\cftsubsubsecfont}{\footnotesize}
\renewcommand{\cftsecpagefont}{\footnotesize}
\renewcommand{\cftsubsecpagefont}{\footnotesize}
\renewcommand{\cftsubsubsecpagefont}{\footnotesize}
\usepackage{etoc}

% --- 8. Abschnittsüberschriften ---
\usepackage{titlesec}
\titleformat{\section}{\normalfont\large\bfseries}{\thesection}{1em}{}
\titleformat{\subsection}{\normalfont\normalsize\bfseries}{\thesubsection}{0.5em}{}
\titleformat{\subsubsection}{\normalfont\normalsize\bfseries}{\thesubsubsection}{0.5em}{}
\setcounter{secnumdepth}{4}

% --- 9. Abstrakt-Umgebung ---
\usepackage{changepage}
\renewenvironment{abstract}
  {
    \begin{adjustwidth}{1.5cm}{1.5cm}
    \small
    \textsc{Abstract. –}%
  }
  {
    \end{adjustwidth}
  }

% --- 10. Todos/Notizen ---
\usepackage{todonotes}

% --- 11. Rahmen/Box-Umgebungen ---
\usepackage{mdframed}
\usepackage{tcolorbox}
\colorlet{shadecolor}{gray!25}

\newenvironment{customTheorem}
  {\vspace{10pt}%
   \begin{mdframed}[
     backgroundcolor=gray!20,
     linewidth=0pt,
     innertopmargin=10pt,
     innerbottommargin=10pt,
     skipabove=\dimexpr\topsep+\ht\strutbox\relax,
     skipbelow=\topsep,
   ]}
  {\end{mdframed}
   \vspace{10pt}%
  }

% --- 12. Python-Integration ---
% (Deaktiviert in dieser Version, aktiviere bei Bedarf)
% \usepackage{pythontex}
% \usepackage[makestderr]{pythontex}

% --- 13. Literaturverwaltung ---
\usepackage{csquotes}
\usepackage[backend=biber, style=alphabetic, citestyle=alphabetic]{biblatex}
\addbibresource{bibliography.bib}

% --- 14. Hyperlinks ---
\usepackage{hyperref}
\hypersetup{
  colorlinks   = true,
  urlcolor     = blue,
  linkcolor    = blue,
  citecolor    = blue,
  frenchlinks  = true
}

% --- 15. Absatzeinstellungen ---
\usepackage[parfill]{parskip}
\sloppy

% --- 16. Umgebungen ---
\usepackage{thmtools}

\newcommand{\CustomHeading}[3]{%
  \par\medskip\noindent%
  \textbf{#1 #2} \textnormal{(#3)}.\enskip%
}

\newenvironment{DEF}[2]{\begin{unitbox}\CustomHeading{Definition}{#1}{#2}}{\end{unitbox}}
\newenvironment{PROP}[2]{\begin{unitbox}\CustomHeading{Proposition}{#1}{#2}}{\end{unitbox}}
\newenvironment{THEO}[2]{\begin{unitbox}\CustomHeading{Theorem}{#1}{#2}}{\end{unitbox}}
\newenvironment{LEM}[2]{\begin{unitbox}\CustomHeading{Lemma}{#1}{#2}}{\end{unitbox}}
\newenvironment{KORO}[2]{\begin{unitbox}\CustomHeading{Corollar}{#1}{#2}}{\end{unitbox}}
\newenvironment{REM}[2]{\begin{unitbox}\CustomHeading{Remark}{#1}{#2}}{\end{unitbox}}
\newenvironment{EXA}[2]{\begin{unitbox}\CustomHeading{Example}{#1}{#2}}{\end{unitbox}}
\newenvironment{STUD}[2]{\begin{unitbox}\CustomHeading{Study}{#1}{#2}}{\end{unitbox}}
\newenvironment{CONC}[2]{\begin{unitbox}\CustomHeading{Concept}{#1}{#2}}{\end{unitbox}}
\newenvironment{OTH}[2]{\begin{unitbox}\CustomHeading{Other}{#1}{#2}}{\end{unitbox}}
\newenvironment{EXE}[2]{\begin{unitbox}\CustomHeading{Exercise}{#1}{#2}}{\end{unitbox}}
\newenvironment{MOT}[2]{\begin{unitbox}\CustomHeading{Motivation}{#1}{#2}}{\end{unitbox}}
\newenvironment{PROOF}[2]{\begin{unitbox}\CustomHeading{Proof}{#1}{#2}}{\end{unitbox}}

% --- Unit Umgebung für Source-Inhalte ---
\usepackage{mdframed}
\newmdenv[
  linewidth=1pt,
  topline=false,
  bottomline=false,
  rightline=false,
  leftmargin=0cm,
  rightmargin=0cm,
  skipabove=10pt,
  skipbelow=10pt,
  innertopmargin=0.5\baselineskip,
  innerbottommargin=0.5\baselineskip,
  backgroundcolor=gray!10,
  linecolor=gray
]{unitbox}

\newenvironment{unit}[1]
  {\begin{unitbox}\textbf{Unit #1}\par\smallskip}
  {\end{unitbox}}

% --- 17. ... ---

% --- 18. Extras ---
\usepackage{stmaryrd}
\usepackage{bbold}  % falls du \mathbb{1} nutzen willst

% --- 00. Titel und Autor ---
\title{Mein Titel}
\author{Tim Jaschik}
\date{\today}

\begin{document}

\maketitle
\rule{\textwidth}{0.5pt}
\begin{abstract}
Kurze Beschreibung …
\end{abstract}
\rule{\textwidth}{0.5pt}
\vspace{0.5cm}

\tableofcontents

\pagebreak




\section{Tutorial: Tensor-Calaculations}

\section*{Prolog: Warum Tensorfelder?}

\subsection*{Die 2-Sphäre als Mannigfaltigkeit}
Sei
\[
  S^{2}_{r} \;=\; \bigl\{ (x,y,z)\in\mathbb R^{3}\mid x^{2}+y^{2}+z^{2}=r^{2} \bigr\}.
\]
$S^{2}_{r}$ ist eine kompakte, orientierbare, zweidimensionale, glatte Mannigfaltigkeit
ohne Rand.  Der Tangentialraum in $p\in S^{2}_{r}$ ist
\[
  T_{p}S^{2}_{r} \;=\; \bigl\{\,v\in\mathbb R^{3}\bigm| v\cdot p = 0\bigr\}.
\]

\subsection*{Tensorfelder (global)}
Ein \((k,\ell)\)\mbox{-}Tensorfeld auf $S^{2}_{r}$ ist eine glatte Abbildung
\[
  T\colon p \longmapsto T(p)\in \underbrace{T_{p}S^{2}{}^{\otimes k}}_{k\text{-mal}}
  \otimes \underbrace{T_{p}^{*}S^{2}{}^{\otimes \ell}}_{\ell\text{-mal}} .
\]
Seine Definition kommt **ohne** Koordinaten aus.

\subsection*{Von der globalen Definition zu Komponenten}

Sei $T$ ein $(k,\ell)$–Tensorfeld auf einer $n$–dimensionalen Mannigfaltigkeit $M$.
Wähle eine Karte $\phi\colon U\to\R^{n}$ mit lokalen Koordinaten $(x^{1},\dots,x^{n})$.
Dann bilden die Vektorfelder
\[
  \partial_{i}\;=\;\frac{\partial}{\partial x^{i}}
\qquad(i=1,\dots,n)
\]
sowie die 1-Formen
\[
  dx^{j}\bigl(\partial_{i}\bigr)=\delta^{j}_{i}
\]
punktweise duale Basen von $T_{p}M$ bzw.\ $T_{p}^{*}M$.
Somit kann man
\[
  T(p)=
  T^{i_{1}\dots i_{k}}{}_{j_{1}\dots j_{\ell}}(p)\;
  \partial_{i_{1}}\!\otimes\!\dots\!\otimes\!\partial_{i_{k}}
  \otimes
  dx^{j_{1}}\!\otimes\!\dots\!\otimes\!dx^{j_{\ell}}
\]
schreiben.  
Die Funktionen $T^{i_{1}\dots i_{k}}{}_{j_{1}\dots j_{\ell}}\colon\phi(U)\to\R$
heißen die \emph{Komponenten von $T$ in der Karte $\phi$}.

Wechselt man zu einer Karte $\tilde\phi$, so gelten wegen
\(
\tilde\partial_{a}
=\tfrac{\partial x^{i}}{\partial\tilde x^{a}}\partial_{i}
\)
und
\(
d\tilde x^{b}
=\tfrac{\partial\tilde x^{b}}{\partial x^{j}}dx^{j}
\)
die Transformationsgesetze
\[
  \tilde T^{a_{1}\dots a_{k}}{}_{b_{1}\dots b_{\ell}}
  =
   \frac{\partial \tilde x^{a_{1}}}{\partial x^{i_{1}}}\!\dots
   \frac{\partial \tilde x^{a_{k}}}{\partial x^{i_{k}}}
   \frac{\partial x^{j_{1}}}{\partial \tilde x^{b_{1}}}\!\dots
   \frac{\partial x^{j_{\ell}}}{\partial \tilde x^{b_{\ell}}}\,
   T^{i_{1}\dots i_{k}}{}_{j_{1}\dots j_{\ell}}.
\]
Damit ist $T$ unabhängig von der Koordinatenwahl.



\subsection*{Zwei Karten und der Metriktensor}

\paragraph{Karte (I): Kugelkoordinaten.}
\[
\Phi(\theta,\varphi)=\bigl(\sin\theta\cos\varphi,
                           \sin\theta\sin\varphi,
                           \cos\theta\bigr),\quad
0<\theta<\pi.
\]
Lokaler Rahmen: $\partial_\theta,\partial_\varphi$;
Korahmen: $d\theta,d\varphi$.

\paragraph{Karte (II): stereographisch (Nordpol).}
\[
\Psi(u,v)=\left(\frac{2u}{1+\rho^{2}},\;
                \frac{2v}{1+\rho^{2}},\;
                \frac{\rho^{2}-1}{1+\rho^{2}}\right),
\qquad \rho^{2}=u^{2}+v^{2}.
\]
Lokaler Rahmen: $\partial_{u},\partial_{v}$;
Korahmen: $du,dv$.

\paragraph{Metrik-Komponenten.}
\[
\begin{aligned}
g_{\theta\theta}&=1,& g_{\theta\varphi}&=0,&
g_{\varphi\varphi}&=\sin^{2}\theta,\\[4pt]
g_{uu}&=g_{vv}=\frac{4}{(1+\rho^{2})^{2}},&
g_{uv}&=0.
\end{aligned}
\]

\paragraph{Transformationstest.}
Mit
$\partial_u\theta=\dfrac{2u}{\rho(1+\rho^{2})}$,
$\partial_u\varphi=-\dfrac{v}{\rho^{2}}$
und
$\sin\theta=\dfrac{2\rho}{1+\rho^{2}}$
folgt
\[
g_{uu}
=(\partial_u\theta)^{2}
  +\sin^{2}\theta\,(\partial_u\varphi)^{2}
=\frac{4}{(1+\rho^{2})^{2}}.
\]
Die übrigen Komponenten ergeben sich analog – die
Transformationsregel ist erfüllt.






\pagebreak

% UNIT 1D — Vom abstrakten Tensorfeld zu Komponenten

CONCEPT: (r,s)-Tensorfeld als multilineare Abbildung  
Ein (r,s)-Tensorfeld \(T\) ordnet jedem Punkt \(p\in M\) eine multilineare
Abbildung  
\[
T_p : \underbrace{T_p^*M \times \dotsm \times T_p^*M}_{r\text{-mal}}
      \times
      \underbrace{T_p M \times \dotsm \times T_p M}_{s\text{-mal}}
      \;\longrightarrow\; \mathbb{R}
\]
zu.  Abstrakte Indizes \(T^{ab}{}_{cd}\) identifizieren dabei lediglich
die Slots (erste zwei „oben“ = kontravariant, letzte zwei „unten“ =
kovariant); sie sind \emph{keine Zahlen}.

DEFINITION: Komponentendarstellung \(T^{\mu\nu}{}_{\kappa\lambda}\)  
Wählt man eine lokale Kartenbasis \(\{\partial_\mu\}\) für \(TM\) und
die duale Basis \(\{dx^\mu\}\) für \(T^*M\), so definiert man  
\[
T^{\mu\nu}{}_{\kappa\lambda}(x)
\;:=\;
T_x\!\bigl(dx^\mu,dx^\nu,\partial_\kappa,\partial_\lambda\bigr).
\]
Damit wird jeder Slot durch Einsetzen eines Basis- (bzw. Dualbasis-)
vektors „gesättigt“ – das liefert eine \emph{Zahl} und damit die
Komponente des Tensors an der Stelle \(x\).

EXAMPLE: Bedeutung von \(T^{ij}{}_{kl}\)  
Schreibt man \(T^{ij}{}_{kl}\) ohne konkrete Zahlenwerte, meint man:

• \(i,j\) belegen die beiden kontravarianten Slots (Vektorargumente),  

• \(k,l\) belegen die beiden kovarianten Slots (1-Form-Argumente).  
Ersetzt man \(i,j,k,l\) durch Karte-Indizes \(\mu,\nu,\kappa,\lambda\),
erhält man die tatsächlich rechenbare Komponente
\(T^{\mu\nu}{}_{\kappa\lambda}\).

PROPOSITION: Transformationsgesetz der Komponenten  
Unter einem Kartenwechsel \(x^\mu \mapsto x^{\mu'}\) transformieren die
Komponenten gemäß  
\[
T^{\mu'\nu'}{}_{\kappa'\lambda'}
=
\frac{\partial x^{\mu'}}{\partial x^{\mu}}
\frac{\partial x^{\nu'}}{\partial x^{\nu}}
\frac{\partial x^{\kappa}}{\partial x^{\kappa'}}
\frac{\partial x^{\lambda}}{\partial x^{\lambda'}}
\,
T^{\mu\nu}{}_{\kappa\lambda},
\]
was genau die Bedingung ist, dass das gesamte Objekt
kartenunabhängig (also ein Tensorfeld) bleibt.



---------------------------------------------------------------------

% UNIT 1E — Indexposition und Rollen der Metrik

CONCEPT: Slot-Typ Indexposition  

• Ein \emph{oben} gesetzter Index steht für einen Vektor-Slot
  (Element aus \(TM\)).  

• Ein \emph{unten} gesetzter Index steht für einen 1-Form-Slot
  (Element aus \(T^*M\)).  
Die Position zeigt also, welchen „Argumenttyp“ der Slot verlangt.

DEFINITION: Heben/Senken mittels Metrik  
Existiert eine Riemannsche Metrik \(g\), so liefert  
\(g^\sharp:T^*M\!\to\!TM,\;\alpha\mapsto g^{\sharp}(\alpha)\)
und \(g^\flat:TM\!\to\!T^*M,\;v\mapsto g^\flat(v)\).
Damit kann man Komponenten umschreiben, z.\,B.  
\(v_i := g_{ij}v^j\)  und  \(T_{ij} := g_{ip}g_{jq}T^{pq}\).

EXAMPLE: Ricci-Tensor als Kontraktion  
Der Riemann-Tensor \(R_{ijk\ell}\) hat Typ (0,4).  Mit der Metrik hebt
man zwei Indizes:  
\(R^i{}_{k} := g^{j\ell}R_{jk\ell}{}^{i}\).  Diese gemischten
Komponenten bilden den Ricci-Tensor, Typ (1,1).

PROPOSITION: Slot-Konstanz beim Heben/Senken  
Heben oder Senken ändert \emph{nicht} die Anzahl der Slots, sondern nur
deren Typ. Die algebraischen Symmetrien des Ursprungstensors bleiben
erhalten, da \(g_{ij}=g_{(ij)}\) symmetrisch ist.



% UNIT 1F — Transformationsverhalten kontravariant vs. kovariant

CONCEPT: Jacobimatrix und ihre Inverse  
Sei \(x^\mu\) ein Koordinatensystem, \(x^{\mu'}=x^{\mu'}(x^1,\dots,x^n)\)
ein neues.  Die \emph{Jacobimatrix} ist  
\(J^{\mu'}{}_{\mu}:=\partial x^{\mu'}/\partial x^{\mu}\);  
ihre Inverse nennen wir \(J^{\mu}{}_{\mu'}\).

DEFINITION: Kontravariantes Transformationsgesetz  
Die Komponenten eines Vektorfeldes \(V^ \mu\) transformieren mit der
Jacobimatrix:
\[
V^{\mu'} \;=\; J^{\mu'}{}_{\mu}\,V^{\mu}.
\]

DEFINITION: Kovariantes Transformationsgesetz  
Die Komponenten einer 1-Form \(\omega_\mu\) benutzen die inverse
(transponierte) Jacobimatrix:
\[
\omega_{\mu'} \;=\; J^{\mu}{}_{\mu'}\,\omega_{\mu}.
\]
Damit gilt \(\omega_{\mu'}V^{\mu'}=\omega_{\mu}V^{\mu}\), d.\,h.\ das
natürliche Paarungsprodukt ist koordinateninvariant.

EXAMPLE: 2-D Kartentransformation (kartesisch → polar)  
\(x^1=r\cos\varphi,\; x^2=r\sin\varphi\).  
Jacobimatrix  
\(J^{\mu'}{}_{\mu}= \begin{pmatrix}\cos\varphi & \sin\varphi \\ -\tfrac{\sin\varphi}{r} & \tfrac{\cos\varphi}{r}\end{pmatrix}\).  
Für einen rein radialen Vektor \((V^x,V^y)=(1,0)\) erhält man  
\(V^{r}=1,\; V^{\varphi}=0\);  
für die Gradienten-1-Form \(d r = (\cos\varphi,\sin\varphi)\) ergibt
sich \((d r)_{r}=1,\;(d r)_{\varphi}=0\).
Dies illustriert: kontravariantes vs. kovariantes Gesetz benutzen
Jacobi bzw. deren Inverse.

PROPOSITION: Allgemeines Transformationsgesetz für (r,s)-Tensoren  
\[
T^{\mu'_1\dots\mu'_r}{}_{\nu'_1\dots\nu'_s}
=
J^{\mu'_1}{}_{\mu_1}\cdots J^{\mu'_r}{}_{\mu_r}
J^{\nu_1}{}_{\nu'_1}\cdots J^{\nu_s}{}_{\nu'_s}
\,
T^{\mu_1\dots\mu_r}{}_{\nu_1\dots\nu_s}.
\]
Kontravariante Indizes nehmen die Jacobis (\(\partial x'/\partial x\)),
kovariante Indizes die inversen Jacobis.

PROOF: Ableitung aus Push-/Pullback  
• Vektorfelder sind Pushforwards von Kurven und nutzen die direkte
Ableitung \(dx^{\mu'}/dx^{\mu}\).  
• 1-Formen sind Pullbacks von Funktionendifferentialen und ziehen via
Kettenregel mit der inversen Matrix zurück.  
Ersetzt man alle Slots eines Tensors entsprechend, erhält man die oben
angegebene Regel.




% UNIT 1A — Abstrakte vs. Koordinaten-Indizes

CONCEPT: Zwei Ebenen der Schreibweise  
Eine \emph{abstrakte} Indexnotation (z.\,B.\ $T^{ab}{}_c$) verwendet
Indizes ausschließlich als Platzhalter für Tensor-Slots; sie tragen
keine numerischen Werte.  
Die \emph{koordinatenbehaftete} Schreibweise
$\bigl(T^{\mu\nu}{}_{\rho}\bigr)(x)$ fixiert eine Karte $(x^\mu)$ und
liefert konkrete Komponentenfunktionen.





DEFINITION: Abstrakter Index  
Ein Symbol wie $a,b,c,\dots$ ohne Bezug auf ein Koordinatensystem heißt
\emph{abstrakter Index}. Gleiche Buchstaben an verschiedenen Stellen
bezeichnen unterschiedliche Slots; es findet \emph{keine} Summation
statt.

EXAMPLE: Koordinatorfreie vs. Komponentenform  
Für $v\in\Gamma(TM)$ und $\alpha\in\Gamma(T^*M)$ gilt  
\[
\alpha_a v^a \quad\text{(abstrakt)}, \qquad
\alpha_\mu v^\mu \;=\;\sum_{\mu=1}^n \alpha_\mu v^\mu
\quad\text{(in einer Karte $(x^\mu)$).}
\]

PROPOSITION: Basisunabhängigkeit  
Tensorgleichungen in abstrakter Indexnotation sind \textbf{basis-}
unabhängig. Ersetzt man jeden abstrakten Index durch einen konkreten
Koordinatenindex, erhält man automatisch eine Komponentengleichung,
die in jeder Karte gilt.

---------------------------------------------------------------------

% UNIT 1B — Summationskonvention & freie/gebundene Indizes

DEFINITION: Einstein-Summationskonvention  
In Komponentenform wird über jedes Indexpaar, das einmal \emph{unten}
und einmal \emph{oben} erscheint, stillschweigend von $1$ bis
$n=\dim M$ summiert:  
\[
T^\mu{}_\mu \;=\;\sum_{\mu=1}^n T^\mu{}_\mu .
\]

CONCEPT: Frei vs. gebunden  
• \emph{Gebundene} (dummy) Indizes erscheinen genau zweimal, einmal
unten, einmal oben – sie werden summiert und dürfen beliebig
umbenannt werden.  
• \emph{Freie} Indizes bestimmen den Typ des resultierenden Tensors und
müssen auf beiden Seiten einer Gleichung identisch auftreten.

EXAMPLE: Typische Fehlerquelle  
Im Ausdruck $A^\mu B_\nu$ sind $\mu$ und $\nu$ frei. Vertauscht man
versehentlich den Index zu $A^\mu B^\nu$, verschwindet die Summation
über $\nu$: Der Ausdruck ändert Typ und Bedeutung.

---------------------------------------------------------------------

% UNIT 1C — Signatur, (Anti)Symmetrisierung

DEFINITION: Signatur der Metrik  
Sei $(M,g)$ pseudo-Riemannsch. Die \emph{Signatur} von $g$ ist
$(s,n\!-\!s)$, wenn lokal die Diagonalform
$\operatorname{diag}(-1,\dots,-1,\;1,\dots,1)$ mit $s$ negativen
Einträgen existiert. Für reine Riemannsche Geometrie gilt meist
$(0,n)$.

DEFINITION: Symmetrisierung / Antisymmetrisierung  
Für einen Tensor $T_{i_1\dots i_k}$:  
\[
T_{(i_1\dots i_k)}
=\frac{1}{k!}\sum_{\sigma\in S_k}T_{i_{\sigma(1)}\dots i_{\sigma(k)}},
\qquad
T_{[i_1\dots i_k]}
=\frac{1}{k!}\sum_{\sigma\in S_k}
\operatorname{sgn}(\sigma)\,
T_{i_{\sigma(1)}\dots i_{\sigma(k)}}.
\]

EXAMPLE: Krümmungssymmetrien in Kurzschreibweise  
\[
R_{ijkl}=R_{[ij][kl]},\qquad
R_{[ijk]l}=0
\quad\Rightarrow\quad
R_{ijkl}+R_{jkil}+R_{kijl}=0.
\]

PROOF: Kontraktion mit Metrik erhält Symmetrie  
Da $g_{pq}=g_{(pq)}$ symmetrisch ist, bleibt eine (Anti)Symmetrie beim
Heben/Senken erhalten: $R^{p}{}_{qrs}=g^{pt}R_{tqrs}$ besitzt exakt
dieselben Schalt-Symmetrien wie $R_{tqrs}$.





\pagebreak
\printbibliography
\end{document}